\section{Верификация и численная точность}
\label{sec:verification}

\draft{Метод изготовленных решений~\cite{salari2000}: задаётся гладкое
поле $T_*$, из него вычисляются источник и граничные данные.
Для сеточных проверок используется ошибка
$E_T=(\sum_P|P|(T_P-T_*(x_P,z_P))^2/\sum_P|P|)^{1/2}$.
Рисунок~\ref{fig:spatial} показывает пространственную сходимость
и перенос между сетками: последние оценки порядка --- 2,14 для MPFA,
1,89 для NMPFA и 1,98 для переноса.}
\reportfigure{05_spatial_convergence.png}
  {Ошибка температуры на гладком изготовленном решении и при переносе поля
  между сетками. Пунктиром показано квадратичное уменьшение ошибки.}
  {fig:spatial}

\draft{На рис.~\ref{fig:contacts} отдельно проверены горизонтальный
и наклонный контакты с контрастом теплопроводности 100:1 и 1000:1.
Ошибка горизонтального теста достигает уровня округления; наклонный
контакт требует сгущения. Максимальный относительный дефект баланса
в серии --- $6{,}8\cdot10^{-14}$.}
\reportfigure{06_material_contacts.png}
  {Проверка сопряжения материалов: ошибка T, нормированная на перепад
  $\Delta T=20\,{}^\circ$C, при сгущении сетки.}
  {fig:contacts}

\draft{На рис.~\ref{fig:time} пока приведены две вспомогательные проверки
на одной ячейке и расчёт гладкого затухающего поля на неподвижной сетке.
Первые две кривые характеризуют временной интегратор, третья --- решение
нестационарного уравнения теплопроводности на мелкой фиксированной сетке.
Для окончательного отчёта заменить этот рисунок проверкой уравнения
теплопереноса в одном неподвижном и одном деформирующемся слое
с заданными точным решением, граничными и начальными условиями.
Новая проверка деформирующегося слоя ещё не выполнена.}
\reportfigure{07_temporal_convergence.png}
  {Предварительные проверки временной точности; рисунок подлежит замене.
  Контрольная линия соответствует первому порядку.}
  {fig:time}

\draft{Влияние сжатых ячеек проверено через тепловое сопротивление
прослойки мощностью $\delta h=N\varepsilon$:
$\Delta T=q\delta h(1/\lambda_{\mathrm{new}}-1/\lambda_{\mathrm{old}})$
ниже прослойки. При $N=40$, $\varepsilon=0{,}05$ м в стационарном тесте
получено не более $0{,}075\,{}^\circ$C и 0,140\% относительного отличия Ro.
Рисунок~\ref{fig:minimum} сопоставляет эту оценку с чувствительностью полного расчёта.}
\reportfigure{08_minimum_cell_height.png}
  {Влияние суммарной мощности сжатых ячеек на T и Ro в стационарном тесте;
  справа --- изменение результата полного расчёта при другом $\varepsilon$.}
  {fig:minimum}

\draft{На рис.~\ref{fig:easy-accuracy} заданные истории нагрева и охлаждения
сопоставлены с независимым квадратурным интегрированием, проверенным удвоением
числа интервалов. Максимальная относительная ошибка Ro при табулировании
на этих историях --- $1{,}35\cdot10^{-6}\%$.}
\reportfigure{09_easy_ro_accuracy.png}
  {Четыре температурные истории, рассчитанная зрелость и относительная
  ошибка Ro при использовании таблицы относительно независимой квадратуры.}
  {fig:easy-accuracy}

\draft{Синтетическая обратная задача на рис.~\ref{fig:synthetic}:
известный поток по форме data1, геометрия обратного набора, 550 наблюдений
без добавленного шума. Все наблюдения согласованы с допуском, но ранний
поток отличается до 22,68\%; после продолжения подбора --- до 20,30\%.
Точность наблюдений и восстановление всей истории потока --- разные критерии.}
\reportfigure{10_synthetic_inverse.png}
  {Заданный и восстановленные потоки; справа --- распределения относительных
  отклонений синтетических T и Ro до и после продолжения подбора.}
  {fig:synthetic}

\draft{На рабочих обратных данных фиксирован найденный поток и отдельно
изменяются $\Delta z$, $\Delta t$ и $\varepsilon$
(рис.~\ref{fig:sensitivity}). Переход с 1 на 0,25 млн лет меняет T в точках
наблюдений не более чем на $0{,}0013\,{}^\circ$C, Ro --- на 0,120\%.
Это проверка чувствительности, а не оценка ошибки относительно точного решения.}
\reportfigure{11_inverse_discretization_sensitivity.png}
  {Максимальные относительные изменения T и Ro в точках наблюдений
  при изменении одного параметра дискретизации; сравнение с основной сеткой.}
  {fig:sensitivity}
